\section{Conclusion}\label{sec:conclusion}
Our contribution is a filter that turns the shared supply's limits into one exact power budget per substep; with exact battery parameters and under the assumptions of Corollary~\ref{cor:closed}, it guarantees both limits at every substep in which currents meeting the budget exist, for any policy it runs with. Substeps without such currents were rare and led to no violation in our evaluations with exact battery parameters, with a mean infeasible time of at most \res{infeas_exact_max}~s per \res{env_episode_s}~s episode. On a simulated hopping leg whose two drives share one battery and DC bus, every policy run with the filter met the bus-voltage and battery-current limits, whereas every policy that hopped without it breached a limit (Tables~\ref{tab:viol}, \ref{tab:ws} and~\ref{tab:abl}). Policies trained with the filter in the loop did not fall in any of their \res{filt_neps_all} evaluation episodes, whereas policies filtered only at test time fell in some settings: the baseline policies, trained without the filter or a penalty, in up to \res{abl_bft_falls_max} of \res{abl_bft_neps} episodes of a stress-test setting, and one of the \res{n_seeds} seeds of the penalty policies ($w=\res{pst_w}$) in \res{pst_fall_k} of \res{eval_eps_set} episodes in each of \res{pst_nfallset} settings (Table~\ref{tab:abl}). In exchange, the penalty policies with the test-time filter hopped higher in most settings: their mean max height was above that of the filter-trained policies at every state of charge of the main evaluation and in \res{pst_nhigher} of \res{pst_nrows} stress-test combinations of setting and state of charge (Table~\ref{S-tab:wsfull} of the supplementary material and Table~\ref{tab:abl}). When falls are costly, the filter should therefore stay in the loop during fine-tuning; a penalty-trained policy with the filter at test time is an option when a small risk of falling is acceptable for extra height. With exact battery parameters the filter needs no margin: with both margins zero, the peak battery current reached $I_{\lim}=\res{pm_ilim}$~A (\res{pm_deltaV0.0deltaI0.0_filtws_imaxpk}~A at most) without exceeding it (Section~\ref{sec:pjm}). With errors in the battery model the guarantee no longer holds: a plant resistance \res{set_rsmall} times the modelled one made the filter breach the current limit, an extended Kalman filter that estimates the resistance removed these breaches, and measurement noise far above the estimator's assumed sensor noise still led to breaches (Tables~\ref{tab:abl} and~\ref{tab:est}). In practice, the margins therefore need to cover only the errors of the battery estimates.
